\section{nanoGPT: de caracteres a un Decoder-only Transformer capaz de generar}

La explicación anterior basada en grafos todavía puede quedarse en el plano de la intuición. Esta sección sigue el walkthrough de nanoGPT de Karpathy para unir texto, token, batch, block, loss, mask y generation en una ruta ejecutable. El objetivo no es memorizar 300 líneas de código, sino poder responder, para cada dimensión de un tensor, qué representa: qué muestras independientes, qué posiciones temporales, qué heads y qué features.

\subsection{Alcance de la evidencia y Tiny Shakespeare}

Al presentar nanoGPT, Karpathy subraya dos puntos: el código es lo bastante pequeño para que los estudiantes lo lean completo y, a la vez, la implementación es lo bastante rigurosa como para reproducir, en su momento, experimentos de nivel GPT-2 entrenando unas 38 horas con 8 GPU. La demostración de clase usa Tiny Shakespeare y un tokenizer de caracteres para exponer el mecanismo; esto no significa que un GPT de producción use el mismo vocabulary ni la misma configuración de cómputo.

\lecturefigure{slide-32-nanogpt.jpg}{nanoGPT: minimalista, legible y aun así capaz de experimentos serios de reproducción de GPT.}{Video oficial de Stanford Online 00:31:01--00:31:30.}

\teachervoice{Karpathy admite primero que la analogía del grafo «se le ocurrió esta misma mañana» y enseguida pasa a la executable implementation. Vale la pena aprender de este orden: la intuición se encarga de plantear preguntas verificables, y el código y el tensor shape impiden que la intuición se descontrole.}

\lecturefigure{slide-33-tiny-shakespeare-dataset.jpg}{Tiny Shakespeare: el texto de las obras se concatena en una secuencia de caracteres de aproximadamente 1 MB.}{Video oficial de Stanford Online 00:31:31--00:32:00.}

Un modelo de lenguaje no lee el «significado» directamente, sino symbol IDs discretos. Para el conjunto de caracteres $\mathcal V$ se define la correspondencia $\operatorname{encode}:\mathcal V\to\{0,\ldots,|\mathcal V|-1\}$, y el texto completo se convierte en una secuencia unidimensional de enteros. Lo que el modelo aprende es la probabilidad condicional bajo esta codificación y esta distribución de entrenamiento; el «Shakespeare infinito» no es más que muestrear una y otra vez de la learned conditional distribution.

\subsection{Tokenización, límites de documento y secuencia de entrada}

Un \term{token} es la unidad básica discreta que procesa el modelo; puede ser un carácter, un subword, un patch u otro element. El \term{vocabulary} es el conjunto de tokens disponibles. La tokenización por caracteres es sencilla y no tiene unknown words, pero alarga las secuencias; la tokenización por subwords acorta las secuencias, pero sus reglas de codificación son más complejas. El embedding asigna a cada token ID un vector de $d$ dimensiones.

\lecturefigure{slide-34-character-tokenization.jpg}{Tokenizer de caracteres: establece una correspondencia bidireccional entre caracteres e IDs enteros.}{Video oficial de Stanford Online 00:32:00--00:32:31.}

\begin{knowledgebox}{Lectura de la figura: encode/decode no es la inferencia del modelo}
El tokenizer es solo una deterministic interface: \texttt{encode} convierte una cadena en enteros y \texttt{decode} realiza la correspondencia inversa; el modelado probabilístico real ocurre en el embedding, los Transformer blocks y el output head. El valor numérico de un ID de carácter no tiene semántica: el ID 7 no es «mayor» que el ID 3.
\end{knowledgebox}

\lecturefigure{slide-35-tokenized-data.jpg}{Tokenized data: un texto largo se convierte en un stream unidimensional de enteros.}{Video oficial de Stanford Online 00:32:27--00:33:00.}

Varios documentos independientes pueden separarse con un special end-of-text token. En clase se dijo de palabra que el modelo debe aprender a «wipe memory» en el límite; una formulación más precisa es esta: el token de límite le da al modelo una señal de segmentación; si la attention mask todavía permite leer a través del límite, que el modelo ignore el texto previo depende de los datos de entrenamiento, del objetivo y de los pesos aprendidos, no de que el token borre automáticamente algún RNN hidden state.

\begin{lstlisting}[caption={Codificación de caracteres de Tiny Shakespeare y construcción del shifted batch.}]
chars = sorted(list(set(text)))
stoi = {ch: i for i, ch in enumerate(chars)}
encode = lambda s: [stoi[ch] for ch in s]
data = torch.tensor(encode(text), dtype=torch.long)

def get_batch(data, batch_size, block_size):
    starts = torch.randint(len(data) - block_size - 1,
                           (batch_size,))
    x = torch.stack([data[i:i + block_size] for i in starts])
    y = torch.stack([data[i + 1:i + block_size + 1]
                     for i in starts])
    return x, y
\end{lstlisting}

\subsection{Batch size y block size: no confundir las dos dimensiones}

El \term{block size} es la context length máxima de cada secuencia en un forward pass; el \term{batch size} es el número de secuencias independientes que se procesan en paralelo. Si $B=4,T=8$, la entrada $X\in\mathbb N^{4\times8}$ contiene cuatro fragmentos de ocho tokens que no se comunican entre sí. El paralelismo en la GPU ocurre en varias dimensiones (batch, time, head y feature), pero la attention no debe intercambiar información entre batch elements.

\lecturefigure{slide-36-batch-construction.jpg}{Recorte aleatorio de lotes de entrenamiento de $B\times T$ a partir de un token stream unidimensional.}{Video oficial de Stanford Online 00:33:00--00:33:38.}

\begin{knowledgebox}{Lectura de la figura: las cuatro líneas horizontales no son una secuencia más larga}
Cada fila es una muestra independiente, cuyo punto de inicio se muestrea al azar del stream de datos; la longitud de cada fila la controla el block size. Confundir la dimensión batch con la dimensión time provoca fugas de información entre muestras. Aumentar el batch mejora sobre todo el aprovechamiento del hardware y la estimación del gradiente; solo aumentar el block size amplía el contexto que una muestra puede usar directamente, y además incrementa de forma notable el cómputo de attention y la memoria de video.
\end{knowledgebox}

\lecturefigure{slide-37-input-target-tensors.jpg}{Input/target tensors: el objetivo es la entrada desplazada una posición a la derecha en el tiempo.}{Video oficial de Stanford Online 00:33:38--00:35:00.}

Para cada posición $t$, el prefix de entrada es $x_{\le t}$ y la etiqueta de supervisión es $x_{t+1}$. Por ello, un chunk de longitud $T$ aporta a la vez $T$ ejemplos de entrenamiento next-token. Si los logits son $z_{b,t,v}$ y las etiquetas son $y_{b,t}$, la cross-entropy promedio es

\[
\mathcal L=-\frac{1}{BT}\sum_{b=1}^{B}\sum_{t=1}^{T}
\log\frac{\exp z_{b,t,y_{b,t}}}{\sum_{v\in\mathcal V}\exp z_{b,t,v}}.
\]

Aquí, $v$ recorre el vocabulary. En entrenamiento, un solo forward supervisa todas las posiciones temporales, mientras que en generación todavía hay que agregar un token a la vez.

\subsection{Embedding, Block y residual stream}

El token embedding $E_{\text{tok}}[x_t]$ aporta el «qué es» y el position embedding $E_{\text{pos}}[t]$ aporta el «dónde está»:

\[
\mathbf x_t^{(0)}=E_{\text{tok}}[x_t]+E_{\text{pos}}[t].
\]

Estos states entran en blocks repetidos y, al final, pasan por LayerNorm y un language-model head lineal que produce los logits sobre el vocabulary. Decoder-only significa que no hay un encoder independiente ni cross-attention; no significa que el modelo «no tenga decoder blocks».

\lecturefigure{slide-38-decoder-missing-gpt.jpg}{GPT decoder-only: conserva masked self-attention y MLP, y elimina encoder/cross-attention.}{Video oficial de Stanford Online 00:35:00--00:38:31.}

\begin{knowledgebox}{Lectura de la figura: las marcas rojas son una resta de arquitectura}
Cada capa del decoder del Transformer original contiene masked self-attention, cross-attention y feed-forward. GPT no tiene un encoder de secuencia de origen, por lo que elimina la cross-attention; la causal self-attention que queda reúne información del contexto a la izquierda, y el MLP se encarga de la transformación local de cada token. El output head proyecta el hidden state final a los logits del siguiente token.
\end{knowledgebox}

\lecturefigure{slide-39-block-module.jpg}{Block de nanoGPT: dos residual updates, una de attention con pre-norm y otra de MLP.}{Video oficial de Stanford Online 00:37:08--00:38:31.}

La residual connection suma la salida de la subcapa al stream principal, lo que da al gradiente una ruta corta; \term{LayerNorm} normaliza el state de un token individual en la dimensión feature y estabiliza la escala entre capas; \term{pre-norm} significa normalizar antes de entrar en la subcapa. Los parámetros del Block se comparten entre las distintas posiciones de token, de modo que la misma regla de cálculo puede aplicarse a la representación de cualquier posición.

\lecturefigure{slide-40-mlp-module.jpg}{MLP por token: expansión lineal, GELU, proyección lineal y dropout.}{Video oficial de Stanford Online 00:38:31--00:38:55.}

\term{GELU} es una activación no lineal suave. La feed-forward network clásica amplía la residual width de $d$ a aproximadamente $4d$ y luego la proyecta de vuelta a $d$. El MLP no lee información entre tokens, así que su cálculo para todos los tokens puede hacerse en paralelo por completo; sin embargo, las entradas de los distintos tokens ya mezclaron contexto a través de la attention, por lo que el mismo MLP produce salidas distintas en contextos distintos.

\subsection{Causal self-attention: la mask impide la fuga de la respuesta}

Si la posición $t$ pudiera leer posiciones futuras $s>t$, la etiqueta de entrenamiento ya estaría a la derecha en el tensor de entrada y el modelo haría trampa. La causal mask se define como

\[
M_{t,s}=\begin{cases}
0,&s\le t,\\
-\infty,&s>t.
\end{cases}
\]

Al sumarla a los logits antes del softmax, la probabilidad de las posiciones futuras se vuelve cero. La mask forma parte del objetivo de entrenamiento: permite calcular todas las posiciones en paralelo sin dejar de respetar la left-to-right factorization.

\lecturefigure{slide-41-full-implementation.jpg}{Implementación completa de CausalSelfAttention: reshape de Q/K/V, mask, softmax, value aggregation y projection.}{Video oficial de Stanford Online 00:38:55--00:41:34.}

\begin{knowledgebox}{Lectura de la figura: un tensor de alta dimensión es solo el mismo algoritmo procesado por lotes}
La implementación proyecta $X$ y lo reorganiza (reshape) en batch $B$, heads $H$, time $T$ y head feature $d_h$. Una multiplicación de matrices calcula de una sola vez los query-key dot products de todos los batches, heads y posiciones; \texttt{masked\_fill} fija los logits futuros en $-\infty$; el softmax da los weights; y los weights se multiplican por los values. transpose/contiguous/view sobre todo reordenan dimensiones; no son un nuevo mecanismo de aprendizaje.
\end{knowledgebox}

\begin{lstlisting}[caption={Forward mínimo de causal self-attention con anotaciones de forma.}]
def forward(self, x):
    B, T, C = x.shape
    q = self.q(x).view(B, T, self.n_head, C // self.n_head)
    k = self.k(x).view(B, T, self.n_head, C // self.n_head)
    v = self.v(x).view(B, T, self.n_head, C // self.n_head)
    q, k, v = (z.transpose(1, 2) for z in (q, k, v))

    scores = q @ k.transpose(-2, -1) / math.sqrt(k.size(-1))
    scores = scores.masked_fill(self.causal[:, :, :T, :T] == 0,
                                float("-inf"))
    weights = torch.softmax(scores, dim=-1)
    y = weights @ v
    y = y.transpose(1, 2).contiguous().view(B, T, C)
    return self.proj(y)
\end{lstlisting}

\subsection{Curva de entrenamiento y autoregressive generation}

Ya construimos la loss de entrenamiento; esta subsección conecta «la loss baja» con «la generación de texto» y precisa qué puede demostrar cada una. Que la curva de entrenamiento baje indica que el modelo predice mejor los tokens held-out, pero no demuestra por sí solo capacidad gramatical, factual ni de razonamiento. En Tiny Shakespeare, una loss más baja hace que la salida aprenda poco a poco nombres de personajes, saltos de línea, formas de palabras y el formato teatral; como los datos son pocos y el context a nivel de caracteres es limitado, la generación todavía puede ser fluida localmente y carecer de sentido en lo global. Al leer la figura, primero hay que ver si el objetivo de optimización mejora y después revisar si el texto muestreado refleja estadística local, estructura de largo alcance o reproducción memorizada.

\lecturefigure{slide-42-training-loss.jpg}{La loss de entrenamiento baja: la next-character prediction mejora de forma continua.}{Video oficial de Stanford Online 00:41:34--00:41:43.}

\begin{warningbox}{Lectura de la figura: la loss es evidencia sobre la función objetivo, no una escala de inteligencia general}
La cross-entropy mide la log-probabilidad del verdadero token siguiente. Respalda la conclusión de que «el modelo se ajusta mejor a esta distribución de datos», pero por sí sola no demuestra coherencia a largo plazo, exactitud factual, robustez ni seguridad. Además, las curvas de entrenamiento y de validación deben leerse por separado para distinguir memorization de generalization.
\end{warningbox}

\lecturefigure{slide-43-fake-shakespeare-output.jpg}{«Falso Shakespeare» generado por muestreo: el modelo ya absorbió el formato y la estadística local.}{Video oficial de Stanford Online 00:41:41--00:43:12.}

La generación parte de un token inicial; en cada paso toma los logits de la última posición, muestrea un token nuevo, lo agrega al context y repite. Si el context máximo es $T_{\max}$, en el paso $n$ solo se conserva la ventana más reciente:

\[
\hat x_{n+1}\sim p_\theta(\cdot\mid \hat x_{\max(1,n-T_{\max}+1):n}).
\]

\begin{lstlisting}[caption={Autoregressive generation con context cropping limitado.}]
@torch.no_grad()
def generate(model, ids, max_new_tokens, block_size):
    for _ in range(max_new_tokens):
        context = ids[:, -block_size:]
        logits, _ = model(context)
        probs = torch.softmax(logits[:, -1, :], dim=-1)
        next_id = torch.multinomial(probs, num_samples=1)
        ids = torch.cat((ids, next_id), dim=1)
    return ids
\end{lstlisting}

\teachervoice{En clase, al toy model con block size 8 se le «devuelven» paso a paso los caracteres que predice: el context crece de 1 a 8 y, al rebasarlo, la ventana empieza a deslizarse. Esta demostración importa más que solo decir «autorregresivo», porque expone la tensión entre el entrenamiento en paralelo, la generación en serie y el context limitado.}

\subsection{Resumen de la sección}

La ruta de nanoGPT lleva los mecanismos abstractos al nivel del tensor: el texto pasa por el tokenizer y se convierte en un stream de IDs; el recorte aleatorio forma batches de $B\times T$; la entrada y el objetivo se desfasan una posición; los token/position embeddings entran en decoder blocks repetidos; la causal mask impide la fuga del futuro; la cross-entropy entrena todas las posiciones; y la generation muestrea un token a la vez, limitada por el block size.
